%% ma: L12-B17-0014
%% lop: 12
%% chuong: 5
%% bai: 17
%% dang: 12.17.9
%% loai: TLN
%% muc: VDC
%% dap_an: "738"
%% nguon: "(NBV)-10 ĐỀ KIỂM TRA HK2-TOÁN 12-DUNG CHUNG 3 SÁCH.pdf (H01, Đề số 1), Phần 3 Câu 2"
%% kien_thuc: "Bài 17 (mặt cầu, vị trí của điểm đối với mặt cầu), tam giác đồng dạng, bất đẳng thức tam giác"
%% trang_thai: cho-duyet
%% ly_do: "Dạng mới đề xuất 12.17.9 (giá trị nhỏ nhất của MA + kMB với M thuộc mặt cầu, dựng điểm phụ theo tỉ số đồng dạng), chờ giáo viên duyệt."
%% ghi_chu: "Đáp án gốc 738 đúng (giải độc lập: I(1;0;2), R = √10, IA = 2R, C(1;1/2;1/2), MA = 2MC, 2BC = √738). Đề in A(1;2;-4), B(1;2;14)."
\begin{ex}
Trong không gian $Oxyz$, cho mặt cầu $(S)\colon(x-1)^2+y^2+(z-2)^2=10$ và hai điểm $A(1;2;-4)$, $B(1;2;14)$. Một điểm $M$ thay đổi trên $(S)$ sao cho giá trị nhỏ nhất của $MA+2MB$ bằng $\sqrt a$. Tính $a$.
\shortans{738}
\loigiai{Mặt cầu $(S)$ có tâm $I(1;0;2)$, bán kính $R=\sqrt{10}$. Ta có $IA=\sqrt{0+4+36}=2\sqrt{10}=2R$.

Gọi $C$ là điểm thuộc đoạn $IA$ sao cho $\overrightarrow{IC}=\dfrac14\overrightarrow{IA}$, tức $C\left(1;\dfrac12;\dfrac12\right)$. Với $M$ bất kì trên $(S)$: $\dfrac{IC}{IM}=\dfrac{IM}{IA}=\dfrac12$ và góc $\widehat{MIC}=\widehat{AIM}$ nên $\triangle ICM\backsim\triangle IMA$, suy ra $\dfrac{MC}{MA}=\dfrac12$, tức $MA=2MC$.

Vì $IC=\dfrac R2<R$ nên $C$ nằm trong $(S)$, còn $IB=\sqrt{148}>R$ nên $B$ nằm ngoài $(S)$. Do đó $MA+2MB=2(MC+MB)\ge2BC$, dấu bằng khi $M$ là giao điểm của đoạn $BC$ với $(S)$.

$BC=\sqrt{0+\left(\dfrac32\right)^2+\left(\dfrac{27}2\right)^2}=\sqrt{\dfrac{738}4}$ nên $\min(MA+2MB)=2BC=\sqrt{738}$. Vậy $a=738$.}
\end{ex}
